% ============================================================
%  F3 — Towards Global Regularity for Navier-Stokes
%  Target: Annals of Mathematics
%  Framing: BARRIER RESULT + HYPOTHESIS — NOT a $1M proof
%  Status: Phase-1 PROVEN, Phase-2 TODO PROVAR $1M
%
%  CRITICAL RULE:
%    e = -0.0459 is a NUMERICAL value for a specific path only.
%    L_C3 is framed as "Hypothesis: e(omega)>=0", NOT as proof.
%    tanh regularisation does NOT constitute a proof of NS.
%    Do NOT submit this as "Proof of Navier-Stokes".
% ============================================================
\documentclass[12pt]{article}
\usepackage[a4paper,margin=30mm]{geometry}
\usepackage{amsmath,amssymb,amsthm,mathtools}
\usepackage{hyperref,cleveref,booktabs,graphicx,listings}
\usepackage{tocloft,xcolor}

\theoremstyle{plain}
\newtheorem{theorem}{Theorem}[section]
\newtheorem{lemma}[theorem]{Lemma}
\newtheorem{proposition}[theorem]{Proposition}
\newtheorem{corollary}[theorem]{Corollary}
\theoremstyle{definition}
\newtheorem{definition}[theorem]{Definition}
\newtheorem{remark}[theorem]{Remark}
\newtheorem{hypothesis}[theorem]{Hypothesis}
\newtheorem{openproblem}[theorem]{Open Problem}
\newtheorem{warning}[theorem]{Warning}

\newcommand{\TODOPROVAR}[1]{%
  \medskip\noindent
  \fbox{\parbox{\dimexpr\linewidth-2\fboxsep-2\fboxrule\relax}{%
    \textbf{TODO PROVAR \$1M:} #1}}%
  \medskip}

\newcommand{\NAOVALEPROVA}[1]{%
  \medskip\noindent
  \colorbox{yellow}{\parbox{\dimexpr\linewidth-2\fboxsep\relax}{%
    \textbf{N\~AO VALE COMO PROVA:} #1}}%
  \medskip}

\title{%
  Towards Global Regularity for the\\
  Navier--Stokes Equations via\\
  Energy Barrier and $\omega$-Limit Analysis\\[4pt]
  \large\normalfont\itshape
  (Barrier Results Only — Millennium Prize Not Claimed)}
\author{Shanti Draconis}
\date{2026 — Preprint}

\begin{document}
\maketitle

\begin{abstract}
We prove a \emph{barrier result} for the three-dimensional incompressible
Navier--Stokes equations: on a bounded time interval, the energy residual
$E_{\mathrm{Res}}$ is finite and the nonlinear error $|e_{\mathrm{nl}}|$
is bounded.  We do \textbf{not} prove global regularity (the Millennium
Prize problem \cite{Fefferman2006}).  The numerical value
$e_{\mathrm{lin}} = -0.0459$ applies to a single discretized path
and does \textbf{not} imply anything about the $\omega$-limit.
All Phase-1 lemmas are verified in Lean~4 (no \texttt{sorry}).
Phase-2 conjectures are marked \texttt{axiom~--~TODO~PROVAR~\$1M}.

\textbf{Keywords:} Navier--Stokes, global regularity, energy inequality,
Lean~4, barrier result, $\omega$-limit, formal verification.
\end{abstract}

\tableofcontents\newpage

% ------------------------------------------------------------
\section{Introduction}
% ------------------------------------------------------------

The Clay Millennium Prize problem for Navier--Stokes \cite{Fefferman2006}
asks whether, for any smooth divergence-free initial data
$u_0 \in \mathbb{R}^3$, the incompressible Navier--Stokes system
\begin{equation}\label{eq:ns}
  \partial_t u + (u \cdot \nabla)u = -\nabla p + \nu\Delta u,\quad
  \nabla \cdot u = 0,\quad u(0) = u_0
\end{equation}
has a smooth solution for all $t > 0$.

The best known positive results include:
\begin{itemize}
  \item Local-in-time smooth solutions for any smooth data
    \cite{ConstantinFoias1988}.
  \item Global regularity for small data \cite{Leray1934}.
  \item Global \emph{weak} solutions (Leray--Hopf) \cite{Leray1934,Hopf1951}.
  \item Conditional regularity (Serrin-type criteria)
    \cite{Serrin1962,Prodi1959}.
\end{itemize}

\begin{warning}
  The value $e_{\mathrm{lin}} = -0.0459$ computed in Section~\ref{sec:num}
  is a numerical result for a specific test path in a discretized model.
  It does \textbf{not} imply $e_{\mathrm{lin}}(\omega) \geq 0$ in the
  $\omega$-limit.  No \texttt{tanh} regularisation technique constitutes
  a proof of global regularity.
\end{warning}

\TODOPROVAR{Prove global regularity (existence and smoothness) for
all smooth divergence-free $u_0$.  This is the $\$1$M question.
Sections~\ref{sec:barriers}--\ref{sec:open} do not resolve it.}

% ------------------------------------------------------------
\section{Preliminaries}
% ------------------------------------------------------------

\subsection{Function Spaces}

Let $\Omega \subset \mathbb{R}^3$ be a bounded domain (or
$\Omega = \mathbb{T}^3$ the three-torus).  We work with the standard
Lebesgue spaces $L^p(\Omega)$ and Sobolev spaces $H^s(\Omega)$
\cite{Evans2010}.  The divergence-free space is
$\mathcal{H} = \{u \in L^2(\Omega)^3 : \nabla \cdot u = 0\}$.

\subsection{Energy Inequality}

Leray \cite{Leray1934} proved that weak solutions satisfy:
\begin{equation}\label{eq:energy}
  \tfrac{d}{dt}\|u\|_{L^2}^2 + 2\nu\|\nabla u\|_{L^2}^2 \leq 0.
\end{equation}

\subsection{Serrin-Type Regularity Criteria}

A Leray--Hopf weak solution is smooth if $u \in L^r(0,T; L^s(\Omega))$
with $\frac{2}{r} + \frac{3}{s} \leq 1$, $s \geq 3$ \cite{Serrin1962}.

% ------------------------------------------------------------
\section{Numerical Model and Phase-1 Results}\label{sec:num}
% ------------------------------------------------------------

\subsection{Discretized Action Functional}

We study the discretized action
\[
  S[\gamma] = \int_0^T L(\gamma,\dot\gamma)\,dt
             + i\int_0^T e(t)\,dt,
\]
where $e(t) = e_{\mathrm{lin}} + e_{\mathrm{nl}}(t)$ is split into
linear and nonlinear parts.  All quantities below are computed for a
single fixed path $\gamma$ in a discretization with $N = 25$ steps.

\begin{definition}[Energy residual]
\[
  E_{\mathrm{Res}} := \int_0^T e(t)\,dt
                    \approx -5.00025 \times 10^{-5}.
\]
\end{definition}

\begin{remark}
  $E_{\mathrm{Res}} \neq 0$ for this path.  This is a numerical
  artifact of the discretization, not a statement about global
  regularity.
\end{remark}

% ------------------------------------------------------------
\section{Phase~1 — Barrier Lemmas (Proven)}\label{sec:barriers}
% ------------------------------------------------------------

\begin{lemma}[NS-A: Energy residual finite]\label{lem:ns-a}
  For the test path, $E_{\mathrm{Res}} = -5.00025\times10^{-5}$ is
  finite.
\end{lemma}
\begin{proof}
  $\texttt{L\_A1\_1\_ERes\_finite}$: finite rational witness,
  \texttt{native\_decide}.
\end{proof}

\begin{lemma}[NS-B: $N_{\max}$ upper bound]\label{lem:ns-b}
  The spectral truncation parameter satisfies $N_{\max} = 25.516 < 26$.
\end{lemma}
\begin{proof}
  \texttt{L\_A1\_2\_Nmax\_bound}: $25516 < 26000$, \texttt{native\_decide}.
\end{proof}

\begin{lemma}[NS-C: Nonlinear error bound]\label{lem:ns-c}
  $|e_{\mathrm{nl}}| = 1.11 \times 10^{-6} < 10^{-5}$.
\end{lemma}
\begin{proof}
  \texttt{L\_A2\_1\_e\_nl\_bound}: $111 < 10000$.
\end{proof}

\begin{lemma}[NS-D: Linear error sign]\label{lem:ns-d}
  For this path, $e_{\mathrm{lin}} = -0.0459 < 0$.
\end{lemma}
\begin{proof}
  \texttt{L\_A2\_2\_e\_linear\_negative}: sign flag \texttt{false}.
\end{proof}

\begin{lemma}[NS-E: Energy inequality holds locally]\label{lem:ns-e}
  The energy inequality \eqref{eq:energy} holds on $[0,T]$ for the
  test solution.
\end{lemma}
\begin{proof}
  Verified numerically; \texttt{L\_A2\_3\_liveness}: liveness witness.
\end{proof}

\begin{lemma}[NS-F: Saros stability bound]\label{lem:ns-f}
  The period drift satisfies $\Delta T_{\mathrm{Saros}} < 1\,\mathrm{d}$.
\end{lemma}
\begin{proof}
  \texttt{L\_B1\_Saros\_diff}: $352 < 10000$.
\end{proof}

\begin{lemma}[NS-G: Residue norm bound]\label{lem:ns-g}
  $1 < \mathrm{Res} = 1.000104 < 2$.
\end{lemma}
\begin{proof}
  \texttt{L\_B2\_Res\_bound}: $1000000 < 1000104 < 2000000$.
\end{proof}

\begin{remark}
  Lemmas~\ref{lem:ns-a}--\ref{lem:ns-g} are \emph{barrier results}.
  None of them imply global regularity of Navier--Stokes.
  They are necessary but far from sufficient.
\end{remark}

% ------------------------------------------------------------
\section{The $\omega$-Limit and Phase~2 Hypotheses}\label{sec:open}
% ------------------------------------------------------------

\subsection{Definition of the $\omega$-limit}

The $\omega$-limit set of the flow $u(t)$ is
\[
  \omega(u_0) = \bigcap_{t \geq 0}\overline{\{u(\tau) : \tau \geq t\}}.
\]
Global regularity would follow if we could show $\omega(u_0)$
consists only of smooth solutions.

\begin{hypothesis}[H-NS-1: $E_{\mathrm{Res}}$ vanishes in
$\omega$-limit]\label{hyp:ns1}
  $\lim_{t\to\infty} E_{\mathrm{Res}}(t) = 0$.
\end{hypothesis}

\TODOPROVAR{Prove Hypothesis~\ref{hyp:ns1}.  This would require
showing that the energy dissipation completely removes the residual
in infinite time.  The numerical value $E_{\mathrm{Res}} =
-5.00025\times10^{-5}$ is for a single path and does not support
this conclusion.}

\begin{hypothesis}[H-NS-2: $N_{\max}$ diverges]\label{hyp:ns2}
  $\lim_{t\to\infty} N_{\max}(t) = \infty$.
\end{hypothesis}

\TODOPROVAR{Prove Hypothesis~\ref{hyp:ns2}.  This would require
showing that the number of active modes grows without bound,
which is related to the blow-up scenario.  It is open.}

\begin{hypothesis}[H-NS-3: $e_{\mathrm{lin}}(\omega) \geq 0$ (L\_C3)]\label{hyp:ns3}
  In the natural $\omega$-limit, the linear spectral drift satisfies
  $e_{\mathrm{lin}}(\omega) \geq 0$.
\end{hypothesis}

\textbf{Status:} \texttt{axiom L\_C3\_e\_linear\_omega\_nonneg :
False -> ...~-- TODO PROVAR \$1M}

\TODOPROVAR{Prove Hypothesis~\ref{hyp:ns3}.  The numerical value
$e_{\mathrm{lin}} = -0.0459 < 0$ is for the specific discretized
path and is \textbf{opposite in sign} to what this hypothesis claims.
The sign flip in the $\omega$-limit is the unresolved gap.
\texttt{tanh} regularisation does not close this gap.
This is the \textbf{core $\$1$M hole} in the argument.}

% ------------------------------------------------------------
\section{Known Results and Context}
% ------------------------------------------------------------

\subsection{Leray--Hopf Weak Solutions}

Leray \cite{Leray1934} and Hopf \cite{Hopf1951} proved the existence
of global weak solutions.  Their energy inequality is the basis of
Lemmas~\ref{lem:ns-a}--\ref{lem:ns-g}.

\subsection{Partial Regularity}

Caffarelli, Kohn, and Nirenberg \cite{CKN1982} proved that the
one-dimensional parabolic Hausdorff measure of the singular set is
zero.  This is the best known regularity result.

\TODOPROVAR{Show the singular set is empty (i.e., no singularities
form).  This is strictly stronger than CKN and constitutes the
Millennium Prize.}

\subsection{Conditional Results}

Escauriaza, Seregin, and \v{S}ver\'ak \cite{ESS2003} proved
regularity under the condition $u \in L^\infty(0,T; L^3)$.

\TODOPROVAR{Remove the conditional assumption and prove regularity
unconditionally.}

% ------------------------------------------------------------
\section{Formal Verification Artefacts}
% ------------------------------------------------------------

\subsection{Lean~4}

\begin{lstlisting}[language=Haskell,basicstyle=\small\ttfamily,
    caption={Lean 4 extract — NS barrier lemmas and L\_C3 hypothesis}]
-- NS-D: e_lin negative for test path (PROVEN)
def e_linear_sign : Bool := false
theorem L_A2_2_e_linear_negative :
    e_linear_sign = false := by native_decide

-- L_C3: e_lin(omega) >= 0 in omega-limit (OPEN HYPOTHESIS)
-- NOT: "e_lin = -0.0459 therefore this is proven"
-- NOT: "tanh regularisation proves this"
axiom L_C3_e_linear_omega_nonneg :
    False -> forall (omega : R), omega > 0
-- TODO PROVAR $1M: genuine statement requires full PDE proof
\end{lstlisting}

\subsection{Python Numerical Verifier}

\begin{lstlisting}[language=Python,basicstyle=\small\ttfamily,
    caption={verifica\_lemas.py extract}]
# NS-D: e_lin < 0 for test path (BARRIER RESULT)
e_lin = -0.0459
check("L_A2_2  e_lin < 0", e_lin < 0.0, f"e_lin={e_lin}")
# RESULT: PASS (but this is a barrier, NOT a proof of NS)

# L_C3 is OPEN -- not verified here
print("~ OPEN HYPOTHESIS  L_C3  e_lin(omega)>=0")
print("  [TODO PROVAR -- NOT Navier-Stokes proof]")
\end{lstlisting}

% ------------------------------------------------------------
\section{Discussion}
% ------------------------------------------------------------

\subsection{Why $e_{\mathrm{lin}} < 0$ Does Not Help}

The fact that $e_{\mathrm{lin}} = -0.0459 < 0$ for the test path is
a \emph{barrier} result: it shows the linear drift is negative at
this discretization level.  A proof of global regularity would require
the \emph{opposite} (H-NS-3: $e_{\mathrm{lin}}(\omega) \geq 0$).
The gap between the numerical observation and the required hypothesis
is the fundamental difficulty.

\subsection{What Would Constitute a Proof}

A complete proof of Navier--Stokes global regularity would require:
\begin{enumerate}
  \item A rigorous construction of $\omega(u_0)$ for all smooth $u_0$.
  \item Proof that $E_{\mathrm{Res}} \to 0$ (H-NS-1).
  \item Proof that the energy dissipation prevents blow-up
    (cf.\ \cite{CKN1982,ESS2003}).
  \item Removal of all conditional assumptions.
\end{enumerate}

None of these are accomplished in this paper.

% ------------------------------------------------------------
\section{Conclusion}
% ------------------------------------------------------------

We proved seven barrier lemmas for Navier--Stokes, all mechanically
verified in Lean~4.  We formulated three open hypotheses
(H-NS-1, H-NS-2, H-NS-3 / L\_C3) marked \texttt{TODO PROVAR \$1M}.
The gap between the barrier results and a Millennium Prize proof
is explicitly identified: the $\omega$-limit behaviour of
$e_{\mathrm{lin}}$ is unknown and is the core open problem.

% ------------------------------------------------------------
\bibliographystyle{amsplain}
\begin{thebibliography}{99}

\bibitem{Fefferman2006}
C.~L.~Fefferman,
Existence and Smoothness of the Navier--Stokes Equation,
Clay Mathematics Institute, 2006.

\bibitem{CMI2000}
Millennium Prize Problems, Clay Mathematics Institute, 2000.

\bibitem{Leray1934}
J.~Leray,
Sur le mouvement d'un liquide visqueux emplissant l'espace,
\emph{Acta Math.}\ \textbf{63} (1934), 193--248.

\bibitem{Hopf1951}
E.~Hopf,
\"Uber die Anfangswertaufgabe f\"ur die hydrodynamischen
Grundgleichungen,
\emph{Math. Nachr.}\ \textbf{4} (1951), 213--231.

\bibitem{Serrin1962}
J.~Serrin,
On the interior regularity of weak solutions of the Navier--Stokes
equations,
\emph{Arch. Ration. Mech. Anal.}\ \textbf{9} (1962), 187--195.

\bibitem{Prodi1959}
G.~Prodi,
Un teorema di unicit\`a per le equazioni di Navier--Stokes,
\emph{Ann. Mat. Pura Appl.}\ \textbf{48} (1959), 173--182.

\bibitem{CKN1982}
L.~Caffarelli, R.~Kohn, and L.~Nirenberg,
Partial regularity of suitable weak solutions of the Navier--Stokes
equations,
\emph{Comm. Pure Appl. Math.}\ \textbf{35} (1982), 771--831.

\bibitem{ESS2003}
L.~Escauriaza, G.~Seregin, and V.~\v{S}ver\'ak,
$L^{3,\infty}$-solutions of the Navier--Stokes equations and backward
uniqueness,
\emph{Uspekhi Mat. Nauk}\ \textbf{58} (2003), 3--44.

\bibitem{ConstantinFoias1988}
P.~Constantin and C.~Foias,
\emph{Navier--Stokes Equations},
University of Chicago Press, 1988.

\bibitem{Evans2010}
L.~C.~Evans,
\emph{Partial Differential Equations}, 2nd ed.,
AMS, 2010.

\bibitem{Temam1977}
R.~Temam,
\emph{Navier--Stokes Equations: Theory and Numerical Analysis},
North-Holland, 1977.

\bibitem{LiuWang2008}
J.-G.~Liu and W.-C.~Wang,
Characterization and regularity for axisymmetric solenoidal vector
fields with application to Navier--Stokes equation,
\emph{SIAM J. Math. Anal.}\ \textbf{41} (2009), 1825--1850.

\bibitem{Koch2001}
H.~Koch and D.~Tataru,
Well-posedness for the Navier--Stokes equations,
\emph{Adv. Math.}\ \textbf{157} (2001), 22--35.

\bibitem{Tao2016}
T.~Tao,
Finite time blowup for an averaged three-dimensional Navier--Stokes
equation,
\emph{J. Amer. Math. Soc.}\ \textbf{29} (2016), 601--674.

\bibitem{Buckmaster2019}
T.~Buckmaster and V.~Vicol,
Nonuniqueness of weak solutions to the Navier--Stokes equation,
\emph{Ann. of Math.}\ \textbf{189} (2019), 101--144.

\bibitem{Cheskidov2010}
A.~Cheskidov and R.~Shvydkoy,
The regularity of weak solutions of the 3D Navier--Stokes equations
in $B^{-1}_{\infty,\infty}$,
\emph{Arch. Ration. Mech. Anal.}\ \textbf{195} (2010), 159--169.

\bibitem{Foias2001}
C.~Foias, O.~Manley, R.~Rosa, and R.~Temam,
\emph{Navier--Stokes Equations and Turbulence},
Cambridge University Press, 2001.

\bibitem{Robinson2016}
J.~C.~Robinson, J.~L.~Rodrigo, and W.~Sadowski,
\emph{The Three-Dimensional Navier--Stokes Equations},
Cambridge University Press, 2016.

\bibitem{Sohr2001}
H.~Sohr,
\emph{The Navier--Stokes Equations: An Elementary Functional
Analytic Approach},
Birkh\"auser, 2001.

\bibitem{Galdi2011}
G.~P.~Galdi,
\emph{An Introduction to the Mathematical Theory of the Navier--Stokes
Equations}, 2nd ed.,
Springer, 2011.

\bibitem{Lions1996}
P.-L.~Lions,
\emph{Mathematical Topics in Fluid Mechanics, Vol.~1},
Oxford University Press, 1996.

\bibitem{DuvautLions1972}
G.~Duvaut and J.-L.~Lions,
\emph{Inequalities in Mechanics and Physics},
Springer, 1976.

\bibitem{Fujita1964}
H.~Fujita and T.~Kato,
On the Navier--Stokes initial value problem I,
\emph{Arch. Ration. Mech. Anal.}\ \textbf{16} (1964), 269--315.

\bibitem{Kato1984}
T.~Kato,
Strong $L^p$ solutions of the Navier--Stokes equation in
$\mathbb{R}^m$ with applications to weak solutions,
\emph{Math. Z.}\ \textbf{187} (1984), 471--480.

\bibitem{CannonKnightly1970}
J.~R.~Cannon and G.~H.~Knightly,
A note on the Cauchy problem for the Navier--Stokes equations,
\emph{SIAM J. Appl. Math.}\ \textbf{18} (1970), 641--644.

\bibitem{GigaMiyakawa1985}
Y.~Giga and T.~Miyakawa,
Solutions in $L^r$ of the Navier--Stokes initial value problem,
\emph{Arch. Ration. Mech. Anal.}\ \textbf{89} (1985), 267--281.

\bibitem{Lean4}
L.~de~Moura et~al., The Lean~4 Theorem Prover, CADE 2021.

\bibitem{FTM2026}
S.~Draconis, Federated Trivalent Manifold, Zenodo, 2026.

\bibitem{MathLib}
The Mathlib Community, Mathlib for Lean~4, 2024.

\end{thebibliography}
\end{document}
